
\subsubsection{2.1 Execution Feedback for Code LLMs}

CodeRL \citep{le2022coderl} established the RLEF paradigm using binary execution rewards with an actor-critic architecture. PPOCoder \citep{shojaee2023ppocoder} adapted Proximal Policy Optimization for code generation with execution-based rewards. Both approaches trained for multiple epochs and report final accuracy without learning curves or samples-to-threshold metrics.

\subsubsection{2.2 Multi-Granularity Feedback}

RLTF \citep{liu2023rltf} introduced fine-grained feedback incorporating test pass rates and error type information, comparing it to binary rewards. Their results suggested that richer feedback improves performance. The present work builds on the multi-granularity concept but uses an ''information bandwidth'' framing and different infrastructure.

\subsubsection{2.3 Process Supervision}

Process reward models \citep{lightman2023verify} demonstrated that step-by-step supervision outperforms outcome-only supervision for mathematical reasoning. This principle may transfer to code generation via execution feedback, though this requires explicit annotation in the reasoning domain versus free execution signals in code.
